\section*{Chapitre \ref{ch:m2:the-rates-options-market} --- Le marché des options de taux}

\begin{solution}{exo:m2:the-rates-options-market:1}
$95/\sqrt{252} = 5.98$ points de base par jour~; $7\sqrt{252} = 111.1$ points de base par
an.
\end{solution}

\begin{solution}{exo:m2:the-rates-options-market:2}
L'entreprise craint une hausse des taux~: elle achète un \omterm{def:m2:the-rates-options-market:cap}{cap} sur le taux de son prêt, ou une
\omterm{def:m2:the-rates-options-market:swaption}{swaption payeuse} si elle veut le droit de fixer plus tard son taux. Le fonds de pension craint
une baisse des taux~: il achète une \omterm{def:m2:the-rates-options-market:swaption}{swaption receveuse}, ou un \omterm{def:m2:the-rates-options-market:cap}{floor} sur des revenus à taux variable.
\end{solution}

\begin{solution}{exo:m2:the-rates-options-market:3}
Black modélise le taux comme une loi lognormale, toujours positif~; un forward ou
un strike négatif n'a pas de logarithme. La \omterm{def:m2:the-rates-options-market:normal}{volatilité normale} gère tous les signes sans
paramètre libre, se cote en points de base, l'unité dans laquelle les traders raisonnent déjà, et son
équivalent quotidien se compare directement aux mouvements réalisés~; un modèle lognormal décalé
fonctionne aussi, mais ses cotations dépendent du décalage, qui doit être convenu.
\end{solution}

\begin{solution}{exo:m2:the-rates-options-market:4}
Moitié du straddle~: $7.80 \times 0.0095/\sqrt{2\pi}$ du notionnel, soit USD~2.96
millions. Le point mort du straddle se situe à 75.8 points de base de part et d'autre.
\end{solution}

\begin{solution}{exo:m2:the-rates-options-market:5}
15.0~\%, 45.4~\% et 326.7~\%.
\end{solution}

\begin{solution}{exo:m2:the-rates-options-market:6}
Dans le modèle normal, le taux au fixing est symétrique autour du forward, 4~\%,
donc la probabilité et l'ampleur d'un écart de 50 points de base au-dessus sont les mêmes qu'en dessous~:
\omterm{def:m2:the-rates-options-market:cap}{caplet} et floorlet ont des valeurs égales période par période. Dans le modèle de Black, la
distribution du taux est étirée vers la droite, et, à même
volatilité de Black, le \omterm{def:m2:the-rates-options-market:cap}{floor} à 3.5~\% vaudrait moins que le \omterm{def:m2:the-rates-options-market:cap}{cap} à 4.5~\%.
\end{solution}

\begin{solution}{exo:m2:the-rates-options-market:7}
0.359~\%~: le prix normal à la monnaie par unité d'annuité est
$0.009/\sqrt{2\pi} = 0.00359$, et un call de Black vaut moins que le forward,
donc aucune volatilité de Black ne reproduit ce prix en dessous de ce forward.
\end{solution}

\begin{solution}{exo:m2:the-rates-options-market:8}
La volatilité lognormale est un pourcentage du niveau du taux. À 20~\% pour un taux de 5~\%,
le marché anticipe des mouvements d'environ 100 points de base par an~; à 60~\% pour un taux de 1~\%,
environ 60. En points de base, l'unité dans laquelle une position gagne ou perd, la
seconde devise est la plus calme. Comparer les \omterm{def:m2:the-rates-options-market:normal}{volatilités normales}.
\end{solution}

\begin{solution}{pb:m2:the-rates-options-market:1}
\textbf{1.} Une \omterm{def:m2:the-rates-options-market:swaption}{swaption receveuse} 3y$\times$7y, au strike de 4.5~\%~: rembourser l'obligation par anticipation dans
trois ans met fin à ses paiements à 4.5~\% pendant sept ans, ce qui vaut la peine
précisément lorsque le taux à sept ans est inférieur à 4.5~\%.
\textbf{2.} Annuler un swap dans lequel elle paie 4.5~\% revient à en conclure ensuite un
swap dans lequel elle reçoit 4.5~\%~: une \omterm{def:m2:the-rates-options-market:swaption}{swaption receveuse} aux mêmes conditions.
\textbf{3.} 4.00~\%~; annuité 5.336.
\textbf{4.} Dans la monnaie de 50 points de base~: une \omterm{def:m2:the-rates-options-market:swaption}{swaption receveuse} au strike de 4.5~\% sur un
forward à 4~\%.
\textbf{5.} USD~25.87 millions, 517 points de base de notionnel.
\textbf{6.} USD~176\,826 par point de base de \omterm{def:m2:the-rates-options-market:normal}{volatilité normale}.
\textbf{7.} $-\text{USD}~163\,698$ par point de base~: la \omterm{def:m2:the-rates-options-market:swaption}{swaption receveuse} gagne lorsque les taux
baissent.
\textbf{8.} USD~353\,651.
\textbf{9.} USD~73\,740.
\textbf{10.} Environ USD~240 millions ($176\,826/73\,740 \times 100$ millions).
\textbf{11.} Dans les termes du swap~: le dealer lui facture un spread flottant plus faible,
ou lui verse un montant initial, en contrepartie du droit d'annulation.
\textbf{12.} 63.8 points de base par an~: USD~25.87 millions rapportés à une annuité de dix
ans de 8.11 sur USD~500 millions.
\textbf{13.} Lorsque le refinancement sur sept ans lui coûte moins de 4.5~\%, à peu près
lorsque le taux à sept ans plus son spread de crédit est inférieur au coupon~; le dealer
dénoue le swap le même jour, et la banque ne garde ni l'un ni l'autre.
\textbf{14.} Les investisseurs de l'\omterm{def:m2:the-rates-options-market:swaption}{obligation remboursable par anticipation} sont rémunérés pour l'option par un coupon plus
élevé~; la banque vend la même option au dealer et paie un taux variable comme tout
autre émetteur. Si le dealer paie l'option plus cher que ce que les investisseurs
ont facturé, le coût variable de la banque est plus bas que sur une obligation simple
échangée simplement~; elle atteint aussi les investisseurs qui veulent le coupon supplémentaire.
\textbf{15.} La \omterm{def:m2:the-rates-options-market:swaption}{swaption receveuse} s'enfonce davantage dans la monnaie~: sa valeur monte à
USD~34.76 millions, son delta augmente, son vega baisse à USD~156\,048~; le dealer
doit rééquilibrer à la fois sa couverture en delta et sa couverture en vega.
\textbf{16.} Les dealers vendent la volatilité qu'ils acquièrent, ce qui fait baisser la
volatilité des échéances lointaines, surtout aux échéances et aux tenors des calls.
\textbf{17.} Les investisseurs qui ont besoin d'options~: les investisseurs hypothécaires qui couvrent la
convexité négative, les assureurs et les fonds de pension qui couvrent des garanties, et les
spéculateurs qui jugent la volatilité bon marché.
\textbf{18.} Remboursable chaque année, c'est une \omterm{def:m2:the-rates-options-market:swaption}{swaption} bermudéenne~: la banque
exercera à la meilleure date, qui dépend de toute la courbe future~; la valeur et la
couverture exigent un modèle de la dynamique du cube~; celle-ci est un strip de
\omterm{def:m2:the-rates-options-market:swaption}{swaptions} européennes dont les pondérations changent avec le marché.
\textbf{19.} Résultat à retenir~: \emph{la vega du callable} vaut USD~176\,826 par point de
base de \omterm{def:m2:the-rates-options-market:normal}{volatilité normale}, que le dealer couvre en vendant environ USD~240
millions de straddles 3y$\times$7y à la monnaie.
\textbf{20.} Des émetteurs de dette remboursable par anticipation, qui vendent l'option de leurs obligations aux
dealers, qui la revendent.
\end{solution}

\begin{solution}{iq:m2:the-rates-options-market:1}
La volatilité lognormale mesure les variations proportionnellement au niveau du taux~; la \omterm{def:m2:the-rates-options-market:normal}{volatilité normale}
exprime les variations en points de base, indépendamment du niveau. Elles coïncident à la monnaie
à un facteur près du forward. Cela compte lorsque les taux sont faibles ou négatifs,
où les cotations lognormales explosent ou n'existent pas, pour les smiles, que les deux
modèles façonnent différemment, et pour la couverture~: le delta d'une option diffère
d'un modèle à l'autre.
\iqlookfor{unités, le décrochage aux taux bas, et la dépendance au modèle du smile
et du delta.}
\end{solution}

\begin{solution}{iq:m2:the-rates-options-market:2}
Une \omterm{def:m2:the-rates-options-market:swaption}{swaption payeuse} donne le droit d'entrer dans un swap qui paie le taux fixe au strike, valable si
les taux montent~; une \omterm{def:m2:the-rates-options-market:swaption}{swaption receveuse} donne le droit de recevoir le taux fixe. Les emprunteurs qui comptent fixer
leur dette, et les investisseurs qui se couvrent contre la hausse des taux, achètent des \omterm{def:m2:the-rates-options-market:swaption}{swaptions payeuses}~; les investisseurs qui se couvrent
contre la baisse des taux, les investisseurs hypothécaires qui couvrent le remboursement anticipé et les fonds de pension achètent
des \omterm{def:m2:the-rates-options-market:swaption}{swaptions receveuses}~; les émetteurs d'\omterm{def:m2:the-rates-options-market:swaption}{obligations remboursables par anticipation}, par l'intermédiaire de leurs dealers, vendent des \omterm{def:m2:the-rates-options-market:swaption}{swaptions receveuses}.
\iqlookfor{profils de gains, et les acheteurs et vendeurs naturels.}
\end{solution}

\begin{solution}{iq:m2:the-rates-options-market:3}
Avec $F_T = F + \sigma\sqrt T Z$ et $K = F$, le payeur vaut
$A\,\mathbb{E}[\sigma\sqrt T Z^+] = A\sigma\sqrt T\,\mathbb{E}[Z^+] =
A\sigma\sqrt T/\sqrt{2\pi}$, puisque $\mathbb{E}[Z^+] = \varphi(0) = 1/\sqrt{2\pi}$.
Le straddle vaut le double, soit environ $0.8\,A\sigma\sqrt T$.
\iqlookfor{la dérivation en une ligne et la règle empirique du $0.4$.}
\end{solution}

\begin{solution}{iq:m2:the-rates-options-market:4}
Chaque jour, le straddle couvert rapporte environ son gamma multiplié par le carré du
mouvement, divisé par deux, et paie sa décroissance temporelle~; les deux se compensent lorsque le
mouvement du jour égale la volatilité journalière implicite, soit 95 divisé par $\sqrt{252}$, environ 6 points de base. Les jours
à grands mouvements rapportent, les jours à petits mouvements font perdre~; il faut comparer les mouvements
journaliers réalisés, pondérés par le gamma, à six.
\iqlookfor{gamma contre thêta, et le mouvement journalier de seuil de rentabilité.}
\end{solution}

\begin{solution}{iq:m2:the-rates-options-market:5}
Les émetteurs remboursables vendent leurs options d'achat aux dealers au sein de swaps~; les dealers revendent la
volatilité, concentrée sur les longues échéances et les longues maturités de swap, ce qui fait baisser leur volatilité.
Les investisseurs hypothécaires sont vendeurs d'options de remboursement anticipé et se couvrent
en achetant des \omterm{def:m2:the-rates-options-market:swaption}{swaptions}, ou en négociant dynamiquement dans le sens des mouvements, ce qui
fait monter la volatilité réalisée et la volatilité implicite.
\iqlookfor{les deux flux structurels et leurs signes opposés.}
\end{solution}

\begin{solution}{iq:m2:the-rates-options-market:6}
Utiliser une approximation en forme fermée ou rationnelle de la fonction de Bachelier inverse
(précise à la précision machine près, comme publiée à cet effet) plutôt
qu'un algorithme de recherche de racine, vectorisée sur les options~; traiter explicitement les cas limites (prix inférieurs ou égaux à la
valeur intrinsèque, temps nul, ailes profondes). Tester par allers-retours
sur une grille de moneyness et de maturités, contre une référence de bissection lente, sur
des entrées aléatoires et sur les cas limites~; mesurer le débit sur des
lots de type production.
\iqlookfor{un inverse explicite, les cas limites et les tests par allers-retours.}
\end{solution}
